%%%%%%%%%%%%%%%%%%%%%%%%%%%%%%%%%%%%%%%%%%%%%%%%%%%%%%%%%%%%%%%%%%%%%%%%%%%%%%%%%%%%%%%%%%%%%%%%%%%%%%%%%%
%Write by:ShuwenHe
%Date:20230613
%%%%%%%%%%%%%%%%%%%%%%%%%%%%%%%%%%%%%%%%%%%%%%%%%%%%%%%%%%%%%%%%%%%%%%%%%%%%%%%%%%%%%%%%%%%%%%%%%%%%%%%%%%

%%%%%%%%%%%%%%%%%%%%%%%%%%%%%%%%%%%%%%%%%%%%%%%%%%%%%%%%%%%%%%%%%%%%%%%%%%%%%%%%%%%%%%%%%%%%%%%%%%%%%%%%%%
\documentclass[12pt,twiside,a4paper]{ctexbook}
\usepackage[centertags]{amsmath}
\usepackage{amsfonts}
\usepackage{amsthm}
\usepackage{newlfont}
\usepackage{makeidx}
\usepackage{wasysym}
\usepackage{geometry} 
\usepackage{graphics}
\usepackage{slashbox} 
\usepackage{fancyhdr} 
\usepackage[pdftex]{graphicx}
\usepackage{epstopdf}
\usepackage{cite}
\usepackage{listings}
\usepackage{tocbibind}
\usepackage[numbers,sort&compress]{natbib}

\setlength\parskip{\baselineskip}
\setcounter{tocdepth}{8} % 生成目录层级
\setcounter{secnumdepth}{4}
\renewcommand\thesection{\arabic{section}}
\usepackage[pdfstartview=FitH,CJKbookmarks=true,bookmarks,bookmarksnumbered=true,
    colorlinks=true,citecolor=black,linkcolor=black,anchorcolor=green,urlcolor=black]{hyperref}
\usepackage{titlesec}
\titleformat{\chapter}[display]{\normalfont\huge\bfseries\center}{\chaptertitlename}{1pt}{\Huge}
\titleformat{\section}{\normalfont\Large\bfseries}{\thesection}{1em}{}
\titleformat{\subsection}{\normalfont\large\bfseries}{\thesubsection}{1em}{}
\titleformat{\subsubsection}{\normalfont\normalsize\bfseries}{\thesubsubsection}{1em}{}
\titleformat{\paragraph}[runin]{\normalfont\normalsize\bfseries}{\theparagraph}{1em}{}
\titleformat{\subparagraph}[runin]{\normalfont\normalsize\bfseries}{\thesubparagraph}{1em}{}
\titlespacing*{\chapter} {0pt}{10pt}{10pt}
\titlespacing*{\section} {0pt}{0.5ex plus 1ex minus .2ex}{0.3ex plus .2ex}
\titlespacing*{\subsection} {0pt}{0.25ex plus 1ex minus .1ex}{0.5ex plus .1ex}
\titlespacing*{\subsubsection}{0pt}{3.25ex plus 1ex minus .2ex}{1.5ex plus .2ex}
\titlespacing*{\paragraph} {0pt}{3.25ex plus 1ex minus .2ex}{1em}
\titlespacing*{\subparagraph} {\parindent}{3.25ex plus 1ex minus .2ex}{1em}
\numberwithin{chapter}{part}
\geometry{left=2.0cm,right=20mm,top=25mm,bottom=25mm}
\let\cleardoublepage\clearpage
%%%%%%%%%%%%%%%%%%%%%%%%%%%%%%%%%%%%%%%%%%%%%%%%%%%%%%%%%%%%%%%%%%%%%%%%%%%%%%%%%%%%%%%%%%%%%%%%%%%%%%%%%%

%%%%%%%%%%%%%%%%%%%%%%%%%%%%%%%%%%%%%%%%%%%%%%%%%%%%%%%%%%%%%%%%%%%%%%%%%%%%%%%%%%%%%%%%%%%%%%%%%%%%%%%%%%
%mathematics
\usepackage{amssymb}
\usepackage{diagbox}
%%%%%%%%%%%%%%%%%%%%%%%%%%%%%%%%%%%%%%%%%%%%%%%%%%%%%%%%%%%%%%%%%%%%%%%%%%%%%%%%%%%%%%%%%%%%%%%%%%%%%%%%%%

%%%%%%%%%%%%%%%%%%%%%%%%%%%%%%%%%%%%%%%%%%%%%%%%%%%%%%%%%%%%%%%%%%%%%%%%%%%%%%%%%%%%%%%%%%%%%%%%%%%%%%%%%%
%
%%%%%%%%%%%%%%%%%%%%%%%%%%%%%%%%%%%%%%%%%%%%%%%%%%%%%%%%%%%%%%%%%%%%%%%%%%%%%%%%%%%%%%%%%%%%%%%%%%%%%%%%%%

%%%%%%%%%%%%%%%%%%%%%%%%%%%%%%%%%%%%%%%%%%%%%%%%%%%%%%%%%%%%%%%%%%%%%%%%%%%%%%%%%%%%%%%%%%%%%%%%%%%%%%%%%%
%
\usepackage{tipa}
%%%%%%%%%%%%%%%%%%%%%%%%%%%%%%%%%%%%%%%%%%%%%%%%%%%%%%%%%%%%%%%%%%%%%%%%%%%%%%%%%%%%%%%%%%%%%%%%%%%%%%%%%%

%%%%%%%%%%%%%%%%%%%%%%%%%%%%%%%%%%%%%%%%%%%%%%%%%%%%%%%%%%%%%%%%%%%%%%%%%%%%%%%%%%%%%%%%%%%%%%%%%%%%%%%%%%
\begin{document}
%%%%%%%%%%%%%%%%%%%%%%%%%%%%%%%%%%%%%%%%%%%%%%%%%%%%%%%%%%%%%%%%%%%%%%%%%%%%%%%%%%%%%%%%%%%%%%%%%%%%%%%%%%

\author
{
Peking University\\
北京大学\\
ShuwenHe\\
何书文\\
1201220707@pku.edu.cn
}

%%%%%%%%%%%%%%%%%%%%%%%%%%%%%%%%%%%%%%%%%%%%%%%%%%%%%%%%%%%%%%%%%%%%%%%%%%%%%%%%%%%%%%%%%%%%%%%%%%%%%%%%%%
\centerline{
\includegraphics{shuwenhe.png}}
写好一本书：工匠精神！用心打造！夜深写于北京大学图书馆。作者亲自一线带课，所带学生多人保送或考入清华北大，根据多年清华附中、101中学、人大附中、北大附中、十一学校，考试真题分析经验所得。用此书考上心目中名校学生无数！何书文北京大学硕士，资深数学名师、信息学竞赛算法名师，所带学生多名考入人大附中早培、清华附中优才、101 实验班、北大附中实验班等名校。全国中学数学联赛、全国中学数学竞赛的辅导老师，全国NOI、CSP信息学竞赛辅导名师。何书文老师在北京大学学习期间立志从事教育事业，帮学生授业解惑。何书文老师小学期间学习奥数，并多次获奖，为以后的学习与研究打下良好基础。何书文 老师在中学阶段数学、物理均获奖。何书文老师在小学中学期间一直为数学课代表，中小学大学期间担任班长，何书文老师在北京大学被选为科技一苑苑长，组织北大同学积极参与校各项活动，积极参与校学生会工作，何书文老师被北京大学评为优秀入党积极分子.何书文老师经常参加北京大学数学课题的研讨班。何书文 老师是北京大学数学系暑期学校全国选出40 名优秀中青年数学人才之一，参加伦敦国王学院、美国杜克大学、美国纽约大学、加拿大多伦多大学教授组成的学术研讨班，研究PDE(偏微分方程)，量子力学方面的数学课题的研究工作，并获得优异成绩结业。何书文老师作为项目经理用数学建模方法给大型企业开发软件，用数学方法规划提高企业产能协作效率。何书文 老师致力于数学方面的教学与研究工作，所带多名孩子已经被点优才进入清华附中创新班，101 实验班，人大附中早培班，是家长值得信赖的老师。考上学生继续跟随何书文老师学习全国数学联赛，全国数学竞赛系列课程，同时学习NOI、IOI、ACM算法编程竞赛。
%%%%%%%%%%%%%%%%%%%%%%%%%%%%%%%%%%%%%%%%%%%%%%%%%%%%%%%%%%%%%%%%%%%%%%%%%%%%%%%%%%%%%%%%%%%%%%%%%%%%%%%%%%

%%%%%%%%%%%%%%%%%%%%%%%%%%%%%%%%%%%%%%%%%%%%%%%%%%%%%%%%%%%%%%%%%%%%%%%%%%%%%%%%%%%%%%%%%%%%%%%%%%%%%%%%%%
\title{CSP-J}
\maketitle
\tableofcontents % 显示目录
\newpage
\pagestyle{fancy}
%%%%%%%%%%%%%%%%%%%%%%%%%%%%%%%%%%%%%%%%%%%%%%%%%%%%%%%%%%%%%%%%%%%%%%%%%%%%%%%%%%%%%%%%%%%%%%%%%%%%%%%%%%

\lhead{
\includegraphics{shuwenedu.png}}
\rhead{科技特长生升学规划 何校长 电话微信15010729356}
\lfoot{
\includegraphics{pku.png}算法第一人北大何书文}
\rfoot{改变您家孩子命运的老师}
%%%%%%%%%%%%%%%%%%%%%%%%%%%%%%%%%%%%%%%%%%%%%%%%%%%%%%%%%%%%%%%%%%%%%%%%%%%%%%%%%%%%%%%%%%%%%%%%%%%%%%%%%%

%%%%%%%%%%%%%%%%%%%%%%%%%%%%%%%%%%%%%%%%%%%%%%%%%%%%%%%%%%%%%%%%%%%%%%%%%%%%%%%%%%%%%%%%%%%%%%%%%%%%%%%%%%
\chapter{C++基础思维导图}
\includegraphics[scale=0.15]{csp-j.png}

\chapter{程序基本概念}
\section{数据类型}
使用编程语言进行编程时，需要用到各种变量来存储各种信息。变量保留的是它所存储的值的内存位置。这意味着，当您创建一个变量时，就会在内存中保留一些空间。\\
您可能需要存储各种数据类型（比如字符型、宽字符型、整型、浮点型、双浮点型、布尔型等）的信息，操作系统会根据变量的数据类型，来分配内存和决定在保留内存中存储什么。

\subsection{基本的内置类型}
\begin{table}[h]
    \centering
    \caption{几种基本的 C++ 数据类型}
    \label{tab:example}
    \begin{tabular}{|c|c|}
    \hline
    \textbf{类型} & \textbf{关键字}\\
    \hline
    布尔型 & bool\\
    \hline
    字符型 & char\\
    \hline
    整型 & int\\
    \hline
    浮点型 & float\\
    \hline
    双浮点型 & double\\
    \hline
    无类型 & void\\
    \hline
    宽字符型 & wchar\_t\\
    \hline
    \end{tabular}
\end{table}
一些基本类型可以使用一个或多个类型修饰符进行修饰：\\
signed\\
unsigned\\
short\\
long

\subsection{自定义数据类型}
\subsubsection{struct结构体}
结构体是一种自定义数据类型，可以用来存储不同类型的数据成员。结构体可以包含多个成员变量，每个成员变量可以有不同的数据类型。结构体的定义使用关键字struct，并通过花括号括起来列出成员变量。
\begin{lstlisting}[language=C++]
struct Person {
    std::string name;
    int age;
    double height;
};
\end{lstlisting}
上述示例定义了一个名为Person的结构体，包含三个成员变量：name（字符串类型）、age（整数类型）和height（浮点数类型）。

\subsubsection{class类}
类是一种更高级的自定义数据类型，它允许定义数据成员和成员函数，并将它们封装在一个单独的实体中。类可以用于实现面向对象编程的概念，如封装、继承和多态。\\
类的定义使用关键字class，并通过花括号括起来列出成员变量和成员函数。
\begin{lstlisting}[language=C++]
class Circle {
private:
    double radius;

public:
    Circle(double r) {
        radius = r;
    }

    double getArea() {
        return 3.14159 * radius * radius;
    }
};
\end{lstlisting}
上述示例定义了一个名为Circle的类，包含一个私有成员变量radius和两个公有成员函数：构造函数和计算面积的函数getArea()。\\
类提供了更强大的封装能力，可以通过访问修饰符（如private、public和protected）来控制成员的访问权限。此外，类还支持继承和多态等面向对象编程的特性。\\
通过结构体和类，可以根据具体的需求定义自己的数据类型，并在程序中使用它们来组织和处理数据。自定义数据类型可以提高代码的可读性、可维护性和重用性。

\section{变量类型}

\section{变量作用域}
一般来说有三个地方可以定义变量：\\
在函数或一个代码块内部声明的变量，称为局部变量。\\
在函数参数的定义中声明的变量，称为形式参数。\\
在所有函数外部声明的变量，称为全局变量。\\
作用域是程序的一个区域，变量的作用域可以分为以下几种：\\
局部作用域：在函数内部声明的变量具有局部作用域，它们只能在函数内部访问。局部变量在函数每次被调用时被创建，在函数执行完后被销毁。\\
全局作用域：在所有函数和代码块之外声明的变量具有全局作用域，它们可以被程序中的任何函数访问。全局变量在程序开始时被创建，在程序结束时被销毁。\\
块作用域：在代码块内部声明的变量具有块作用域，它们只能在代码块内部访问。块作用域变量在代码块每次被执行时被创建，在代码块执行完后被销毁。\\
类作用域：在类内部声明的变量具有类作用域，它们可以被类的所有成员函数访问。类作用域变量的生命周期与类的生命周期相同。

\section{局部变量}
在函数或一个代码块内部声明的变量，称为局部变量。它们只能被函数内部或者代码块内部的语句使用。下面的实例使用了局部变量：
\begin{lstlisting}[language=C++]
int sum(){
	int a,b,sum; // 局部变量声明
	a = 1,b = 2; // 实际初始化
	sum = a + b;
	cout<<"sum = "<<sum<<endl;
	return 0;
}
\end{lstlisting}

\section{全局变量}
局部变量和全局变量的名称可以相同，但是在函数内，局部变量的值会覆盖全局变量的值。下面是一个实例：
\begin{lstlisting}[language=C++]
// globalVariable 全局变量
int i = 3;
int globalVariable(){
	int i = 5;
	cout<<"i = "<<i<<endl;
	return 0;
}
\end{lstlisting}

\section{块作用域}
块作用域指的是在代码块内部声明的变量：
\begin{lstlisting}[language=C++]
// blockScope块作用域
int blockScope(){
	int i = 1;
	{
		int i = 2; // 块作用域变量
		cout<<"i = "<<i<<endl;
	}
	cout<<"i = "<<i<<endl;
	return 0;
}
\end{lstlisting}
\section{代码介绍}
\begin{lstlisting}[language=C++]
#include <iostream>
using namespace std;
int main() {
	cout<<"hello"<<endl;
	return 0;
}
\end{lstlisting}
C++头文件<iostream>\\
using namespace std; 告诉编译器使用std命名空间。\\
int main()是主函数，程序从这里开始执行。\\
cout<<"Hello World"; 会在屏幕上显示消息 "Hello World"。\\
下一行 return 0; 终止 main( )函数，并向调用进程返回值 0。\\
在 C++ 中，分号是语句结束符。也就是说，每个语句必须以分号结束。它表明一个逻辑实体的结束。

\section{标识符}

\section{关键字}

\chapter{基本数据类型}

\chapter{程序基本语句}

\chapter{基本运算}

\chapter{数学库常用函数}

\section{算术运算符}
\begin{lstlisting}[language=C++]
// 算术运算符
int arithmeticOperator(){
	int a = 5;
	int b = 3;
	int c;
	cout<<"a = "<<a<<endl;
	cout<<"b = "<<b<<endl;
	c = a + b;
	cout<<"c = a + b = "<<c<<endl;
	c = a - b;
	cout<<"c = a - b = "<<c<<endl;
	c = a * b;
	cout<<"c = a * b = "<<c<<endl;
	c = a / b;
	cout<<"c = a / b = "<<c<<endl;
	c = a % b;
	cout<<"c = a % b = "<<c<<endl;
	int d = 7;
	cout<<"d = "<<d<<endl;
	c = d++;
	cout<<"c = d++ = "<<c<<endl;
	c = d--;
	cout<<"c = d-- = "<<c<<endl;
	return 0;
}
\end{lstlisting}

\section{关系运算符}
\begin{lstlisting}[language=C++]
// 关系运算符
int relationalOperator(){
	int a = 5;
	int b = 3;
	cout<<"a = "<<a<<endl;
	cout<<"b = "<<b<<endl;
	int c;
	if (a == b){
		cout<<"a 等于 b"<<endl;
	}else{
		cout<<"a 不等于 b"<<endl;
	}
	if(a < b){
		cout<<"a 小于 b"<<endl;
	}else{
		cout<<"a 不小于 b"<<endl;
	}
	if(a > b){
		cout<<"a 大于 b"<<endl;
	}else{
		cout<<"a 不大于 b"<<endl;
	}
	return 0;
}
\end{lstlisting}

\section{逻辑运算符}
逻辑运算符在C++中用于解决以下问题：\\
1. 条件判断：逻辑运算符允许程序员在条件语句中对多个条件进行组合判断。通过使用逻辑与运算符（\&\&）和逻辑或运算符（||），可以根据多个条件的组合结果来确定程序的执行路径。\\
2. 循环控制：逻辑运算符在循环语句中起到关键作用，例如在while循环或do-while循环中，使用逻辑运算符可以设置多个条件来控制循环的执行和终止条件。\\
3. 布尔逻辑操作：逻辑运算符允许对布尔值进行操作，将多个布尔值进行组合，从而得到新的布尔值。这对于程序中的条件逻辑判断非常有用。\\
通过使用逻辑运算符，程序员可以根据条件的组合结果来进行复杂的判断和控制，从而实现程序的逻辑流程控制和条件判断。这样可以使程序更加灵活和可控，并能够处理多种不同的情况。\\
C++中的逻辑运算符用于对条件表达式进行逻辑运算，通常返回布尔值（true或false）。以下是C++中常用的逻辑运算符：\\
1. 逻辑与运算符（\&\&）：当且仅当两个操作数都为true时，结果为true。否则，结果为false。\\
\begin{lstlisting}[language=C++]
   bool a = true;
   bool b = false;
   bool result = a && b; // 结果为false
\end{lstlisting}
2. 逻辑或运算符（||）：当至少有一个操作数为true时，结果为true。只有当两个操作数都为false时，结果为false。
\begin{lstlisting}[language=C++]
   bool a = true;
   bool b = false;
   bool result = a || b; // 结果为true
\end{lstlisting}
3. 逻辑非运算符（！）：对操作数进行取反操作，如果操作数为true，则结果为false；如果操作数为false，则结果为true。
\begin{lstlisting}[language=C++]
   bool a = true;
   bool result = !a; // 结果为false
\end{lstlisting}
逻辑运算符通常与条件语句（例如if语句和while循环）一起使用，用于控制程序的执行流程和判断条件的满足情况。\\
C++ 中的逻辑运算符包括逻辑与（\&\&）、逻辑或（||）、逻辑非（!）三种。它们的作用是对逻辑表达式进行求值，以判断表达式的真假。\\
当使用逻辑与（\&\&）时，只有当两个操作数都为真（非零）时，整个表达式才为真，否则为假。因此，如果一个操作数为真，另一个操作数为假，整个表达式的结果就是假。\\
同样的道理，当使用逻辑或（||）时，只有当两个操作数都为假（零）时，整个表达式才为假，否则为真。如果一个操作数为假，另一个操作数为真，整个表达式的结果也是真。\\
逻辑非（!）则是将操作数的真假值取反。如果操作数为真，取反后就是假；如果操作数为假，取反后就是真。\\
因此，当使用逻辑运算符时，需要注意操作数的真假值，以便正确地求出整个表达式的值。
\begin{lstlisting}[language=C++]
// 逻辑运算符
int logicalOperator(){
	int a = 3,b = 5,c;
	cout<<"a = "<<a<<endl;
	cout<<"b = "<<b<<endl;
	if (a&&b){
		cout<<"a&&b条件为 true"<<endl;
	}
	if (a || b){
		cout<<"a||b条件为 true"<<endl;
	}
	// 改变a和b的值
	a = 0;
	b = 5;
	if (a && b){
		cout<<"a&&b条件为 true"<<endl;
	}else{
		cout<<"a&&b条件为 false"<<endl;
	}
	if (!(a&&b)){
		cout<<"!(a&&b)条件为 true"<<endl;
	}
	return 0;
}
\end{lstlisting}

\chapter{注释}
\section{单行多行注释}
C++ 支持单行注释和多行注释。注释中的所有字符会被 C++ 编译器忽略。\\
// - 一般用于单行注释。\\
/* ... */ - 一般用于多行注释。

\chapter{常量}
\section{\#define预处理器}
使用 \#define 预处理器定义常量
\begin{lstlisting}[language=C++]
// #define 预处理器定义常量
#define LENGTH 3
#define WIDTH 2

int areaDefine(){
	int area;
	area = LENGTH * WIDTH;
	cout<<"area = "<<area<<endl;
	return 0;
}
\end{lstlisting}

\section{const 关键字}
\begin{lstlisting}[language=C++]
// 使用 const 前缀声明指定类型的常量
int constConstant(){
	const int LENGTH_ = 3;
	const int WIDTH_ = 2;
	int area;
	area = LENGTH_ * WIDTH_;
	cout<<"area = "<<area<<endl; 
	return 0;
}
\end{lstlisting}

\chapter{修饰符类型}
\begin{lstlisting}[language=C++]
int modifier(){
	short int i; // 有符号短整数
	short unsigned int j;
	j = 50000;
	i = j;
	cout<<"j = "<<j<<endl;	
	cout<<"i = "<<i<<endl;	
	return 0;
}
\end{lstlisting}

\chapter{控制语句}
\section{for循环}
// 高斯求和公式求和1+2+3+...+100
\begin{lstlisting}[language=C++]
#include <iostream>
using namespace std;
int main() {
    int sum = 0;
    sum = (1+100)*100/2;
    cout << "gauss sum=" << sum << endl;
    return 0;
}
\end{lstlisting}

\section{}
for循环求和1+2+3+...+100\\
\begin{lstlisting}[language=C++]
int forSum(){
    int sum = 0;
    for (int i = 1; i <= 100; i++) {
        sum += i;
    }
    cout << "for sum=" << sum << endl;
    return 0;
}
\end{lstlisting}

\chapter{函数}
\section{CSP-J（普及组）2022年T1乘方(pow)}

\chapter{数组}

\chapter{STL}

\chapter{class类}
在C++中，`class`关键字用于定义一个类。类是一种用户自定义的数据类型，用于封装数据和操作。类可以包含成员变量（属性）和成员函数（方法），
用于描述对象的状态和行为。

以下是一个简单的C++类的示例：

\begin{lstlisting}[language=C++]
#include <iostream>
using namespace std;

class Rectangle {
private:
    int length;
    int width;

public:
    // 构造函数
    Rectangle(int l, int w) : length(l), width(w) {}

    // 成员函数
    int getArea() {
        return length * width;
    }

    // 成员函数
    int getPerimeter() {
        return 2 * (length + width);
    }
};

int main() {
    // 创建一个 Rectangle 对象
    Rectangle rect(5, 3);

    // 调用对象的成员函数
    int area = rect.getArea();
    int perimeter = rect.getPerimeter();

    // 打印结果
    cout << "Area: " << area << endl;
    cout << "Perimeter: " << perimeter << endl;

    return 0;
}
\end{lstlisting}
在上述示例中，定义了一个名为`Rectangle`的类，它具有私有成员变量`length`和`width`，以及公有成员函数`getArea()`和`getPerimeter()`。构造函数用于初始化对象的数据成员。
在`main()`函数中，创建了一个`Rectangle`对象`rect`，并使用对象的成员函数`getArea()`和`getPerimeter()`计算矩形的面积和周长。最后，使用`std::cout`打印结果。
类提供了一种组织和封装相关数据和行为的方式，使代码更加模块化和可维护。通过类的实例化，可以创建多个对象，每个对象都有自己的数据和方法。
除了成员变量和成员函数，类还可以包含访问修饰符（如`public`、`private`、`protected`）和其他特性（如继承、多态）等。

\section{构造函数}
在C++中，构造函数是一种特殊的成员函数，用于在创建对象时进行初始化操作。构造函数的名称与类的名称相同，并且没有返回类型（包括`void`）。它可以有参数，
也可以没有参数。\\
构造函数在以下几种情况下会被自动调用：\\
1. 在创建对象时，使用`new`运算符动态分配内存时。\\
2. 在声明对象时，使用类的默认构造函数进行初始化。\\
3. 在将一个对象作为参数传递给函数时，调用拷贝构造函数进行复制。\\
以下是一个简单的示例，展示了如何定义和使用构造函数：
\begin{lstlisting}[language=C++]
#include <iostream>

class MyClass {
private:
    int value;

public:
    // 默认构造函数
    MyClass() {
        value = 0;
        std::cout << "Default constructor called" << std::endl;
    }

    // 带参数的构造函数
    MyClass(int val) {
        value = val;
        std::cout << "Parameterized constructor called" << std::endl;
    }

    // 成员函数
    int getValue() {
        return value;
    }
};

int main() {
    // 使用默认构造函数创建对象
    MyClass obj1;
    std::cout << "Value: " << obj1.getValue() << std::endl;

    // 使用带参数的构造函数创建对象
    MyClass obj2(10);
    std::cout << "Value: " << obj2.getValue() << std::endl;

    return 0;
}
\end{lstlisting}
在上述示例中，定义了一个名为`MyClass`的类，其中包含一个私有成员变量`value`和三个构造函数。默认构造函数用于初始化`value`为0，带参数的构造函数用于将传入的值赋给`value`。\\
在`main()`函数中，首先使用默认构造函数创建了一个`MyClass`对象`obj1`，并通过`getValue()`方法获取对象的值并打印。然后，使用带参数的构造函数创建了另一个对象`obj2`，同样获取并打印了对象的值。\\
构造函数在对象创建时自动调用，用于进行必要的初始化工作。你可以根据需要定义不同的构造函数，以支持不同的初始化方式。\\
希望这个示例对你有帮助！如果你还有其他问题，请随时提问。

\section{拷贝构造函数}
在C++中，拷贝构造函数（Copy Constructor）是一种特殊的构造函数，用于创建一个对象的副本。拷贝构造函数通常以传入对象的引用作为参数，
并使用该对象的数据来初始化新对象。\\
拷贝构造函数在以下情况下会被自动调用：\\
1. 在将一个对象作为参数传递给函数时，进行参数的复制。\\
2. 在使用一个对象初始化另一个对象时，进行对象的复制。\\
3. 在函数返回一个对象时，进行对象的复制。\\
以下是一个简单的示例，展示了如何定义和使用拷贝构造函数：
\begin{lstlisting}[language=C++]
#include <iostream>

class MyClass {
private:
    int value;

public:
    // 默认构造函数
    MyClass() {
        value = 0;
        std::cout << "Default constructor called" << std::endl;
    }

    // 带参数的构造函数
    MyClass(int val) {
        value = val;
        std::cout << "Parameterized constructor called" << std::endl;
    }

    // 拷贝构造函数
    MyClass(const MyClass& other) {
        value = other.value;
        std::cout << "Copy constructor called" << std::endl;
    }

    // 成员函数
    int getValue() {
        return value;
    }
};

void printObject(const MyClass& obj) {
    std::cout << "Object value: " << obj.getValue() << std::endl;
}

int main() {
    // 使用默认构造函数创建对象
    MyClass obj1;
    std::cout << "Value: " << obj1.getValue() << std::endl;

    // 使用带参数的构造函数创建对象
    MyClass obj2(10);
    std::cout << "Value: " << obj2.getValue() << std::endl;

    // 使用拷贝构造函数创建对象的副本
    MyClass obj3 = obj2;
    std::cout << "Value: " << obj3.getValue() << std::endl;

    // 作为函数参数传递对象
    printObject(obj3);

    return 0;
}
\end{lstlisting}
在上述示例中，`MyClass`类定义了默认构造函数、带参数的构造函数和拷贝构造函数。拷贝构造函数以传入对象的引用作为参数，并将传入对象的值复制给新对象的成员变量。\\
在`main()`函数中，首先使用默认构造函数创建了一个`MyClass`对象`obj1`，然后使用带参数的构造函数创建了另一个对象`obj2`。接下来，使用拷贝构造函数将`obj2`复制到新对象`obj3`。最后，通过调用`printObject()`函数将`obj3`作为参数传递给函数。\\
拷贝构造函数在对象的复制过程中起到重要作用，确保新对象与原始对象具有相同的值。如果没有显式定义拷贝构造函数，编译器会自动生成一个默认的拷贝构造函数。\\
需要注意的是，拷贝构造函数的参数通常是`const`引用，以防止在拷

\clearpage
\end{document}
